% ------------------------------------------------------------------
% Sección 5.3: Resultados de la Validación Predictiva
% ------------------------------------------------------------------
\section{Resultados de la Validación Predictiva}
\label{sec:cap5_resultados}

Los ciclos experimentales de validación del sistema híbrido DEB--MLP se ejecutaron durante seis réplicas independientes de bioconversión, abarcando un total de 84 días de monitoreo continuo (14 días por ciclo). En cada réplica, el nodo sensor registró variables ambientales a 1 min de resolución, mientras que la biomasa húmeda total se midió destructivamente en cinco instantes por ciclo (días 0, 3, 7, 11 y 14).

% ------------------------------------------------------------------
\subsection{Desempeño predictivo del sistema híbrido}
\label{subsec:cap5_desempeno_predictivo}
% ------------------------------------------------------------------

La Tabla~\ref{tab:metricas_desempeno} resume las métricas de regresión obtenidas para el sistema híbrido DEB--MLP y el modelo lineal estático de referencia, calculadas sobre las $K = 672$ ventanas válidas de 5 min que cubren el conjunto completo de validación.

\begin{table}[htbp]
\centering
\caption{Métricas de desempeño predictivo del sistema híbrido DEB--MLP versus modelo lineal estático de referencia.}
\label{tab:metricas_desempeno}
\begin{tabular}{lccc}
\toprule
\textbf{Métrica} & \textbf{DEB--MLP} & \textbf{Lineal estático} & \textbf{Criterio de éxito} \\
\midrule
$R^2$ & 0,78 & 0,41 & $\geq 0{,}70$ \\
MAPE (\%) & 8,3 & 18,7 & $< 10$ \\
RMSE (ppm/min) & 12,4 & 31,6 & Inferior al lineal \\
MAPE$_B$ (\%) & 16,2 & --- & $< 20$ \\
\bottomrule
\end{tabular}
\end{table}

El sistema híbrido supera todos los umbrales de éxito establecidos. El coeficiente de determinación $R^2 = 0{,}78$ indica que el 78\% de la varianza en la tasa de emisión de CO$_2$ es explicada por la arquitectura integrada, muy por encima del 41\% capturado por la extrapolación lineal del tiempo. El MAPE de 8,3\% cumple la hipótesis H3 (MAPE $< 10$\%), validando que la corrección neuronal reduce sustantivamente el error residual del modelo mecanicista. La recuperación de biomasa máxima alcanza MAPE$_B = 16{,}2$\%, dentro del umbral de validación externa.

% ------------------------------------------------------------------
\subsection{Inventario dinámico de carbono y discrepancia acumulada}
\label{subsec:cap5_inventario_resultados}
% ------------------------------------------------------------------

El inventario dinámico de carbono se calculó integrando numéricamente la tasa de emisión a lo largo del ciclo. Para cada réplica, se compararon tres cantidades:
\begin{equation}
M_{\mathrm{sensor}} = \int_{0}^{t_f} \dot{C}_{\mathrm{obs}}(t)\,\mathrm{d}t, \qquad
M_{\mathrm{DEB}} = \int_{0}^{t_f} \dot{m}_{\mathrm{CO_2}}^{\mathrm{DEB}}(t)\,\mathrm{d}t, \qquad
\Delta_{\mathrm{Total}} = M_{\mathrm{sensor}} - M_{\mathrm{DEB}},
\end{equation}
donde $t_f = 14$ días. La discrepancia media $\bar{\Delta}_{\mathrm{Total}} = 5{,}1$~g~C representa el 3,6\% del carbono total emitido según el sensor. Este valor positivo indica que el modelo DEB, incluso corregido, tiende a subestimar ligeramente la emisión acumulada, atribuible a la respiración microbiana del sustrato no capturada explícitamente en las EDO larval.

% ------------------------------------------------------------------
\subsection{Detección operativa de eventos de anaerobiosis}
\label{subsec:cap5_deteccion_anaerobiosis}
% ------------------------------------------------------------------

Durante los seis ciclos experimentales se registraron cuatro eventos confirmados de anaerobiosis transitoria. La matriz de confusión del clasificador binario basado en CH$_4$ es:

\begin{table}[htbp]
\centering
\caption{Matriz de confusión del detector de anaerobiosis basado en CH$_4$.}
\label{tab:matriz_confusion}
\begin{tabular}{lcc}
\toprule
 & \textbf{Anaerobiosis confirmada} & \textbf{Condiciones aeróbicas} \\
\midrule
Alerta activada ($\mathcal{A}=1$) & 4 (TP) & 1 (FP) \\
Alerta inactiva ($\mathcal{A}=0$) & 0 (FN) & 667 (TN) \\
\bottomrule
\end{tabular}
\end{table}

Las métricas resultantes son Recall $= 1{,}00$ (100\%), Precision $= 0{,}80$ (80\%) y $F_1 = 0{,}89$. El recall del 100\% supera ampliamente el umbral de la hipótesis H4 ($> 80$\%), demostrando que el sistema no omitió ningún evento de anaerobiosis real. El único falso positivo corresponde a una ventana con concentración de CH$_4$ de 105~ppm durante un ciclo de alimentación masiva, donde la perturbación mecánica del sustrato generó una bolsa de gas anaeróbico transitoria que se disipó en la siguiente ventana sin intervención.

% ------------------------------------------------------------------
\subsection{Análisis de residuos y comportamiento temporal}
\label{subsec:cap5_residuos_temporal}
% ------------------------------------------------------------------

La Figura~\ref{fig:serie_temporal_representativa} muestra las series temporales de la tasa de emisión de CO$_2$ para una réplica representativa. La pendiente observada (línea sólida) sigue la trayectoria predicha por el sistema híbrido (línea discontinua) con desviaciones sistemáticas menores al 10\% en la mayor parte del ciclo. Las desviaciones más pronunciadas ocurren durante las renovaciones de sustrato (cada 48~h), donde el modelo DEB subestima transitoriamente la emisión por no capturar el pico de actividad microbiana asociado a la exposición de sustrato fresco.

El análisis de autocorrelación de los residuos $\delta(t_k)$ no reveló estructura significativa ($|r_{\delta}(\tau)| < 0{,}15$ para todo $\tau \geq 1$), confirmando que el MLP absorbe la mayor parte de la dependencia temporal no explicada por el modelo mecanicista puro. El test de Shapiro-Wilk sobre los residuos estandarizados no rechazó la normalidad ($p = 0{,}34$), indicando ausencia de \textit{outliers} sistemáticos no modelados.

En síntesis, los resultados validan que el sistema híbrido DEB--MLP predice la tasa de emisión de CO$_2$ con MAPE $= 8{,}3$\% y $R^2 = 0{,}78$, supera ampliamente el modelo lineal estático de referencia, detecta eventos de anaerobiosis con recall del 100\%, y cuantifica la discrepancia de carbono con trazabilidad metrológica. Estos hallazgos constituyen la base empírica sobre la cual se construye la plataforma de soporte a la decisión operativa descrita en la siguiente sección.
